\section{Method}
\label{sec:method}

We describe Sparse Feature Preservation (SFP) in three parts: the empirical hypothesis that motivates the method (\cref{sec:hypothesis}), the preservation algorithm (\cref{sec:algorithm}), and the formally verified mathematical properties (\cref{sec:formal}).

\subsection{Hypothesis: Forgetting is Low-Dimensional}
\label{sec:hypothesis}

Let $f_\theta$ be a language model with parameters $\theta$, and let $a_\ell(\theta) \in \R^d$ denote the activation vector at layer $\ell$ for a given input.
When we fine-tune from $\theta^*$ (pre-trained) to $\theta'$ (fine-tuned on a new task), the activation drift at layer $\ell$ is $\Delta a_\ell = a_\ell(\theta') - a_\ell(\theta^*)$.

Let $U \in \R^{d \times k}$ be a PCA basis computed from old-task activations at layer $\ell$, and let $U_r \in \R^{d \times r}$ ($r \ll k \ll d$) be the submatrix corresponding to the top-$r$ most important dimensions, ranked by ablation importance.
We define the \emph{projected drift} in the important subspace as $P_W(\Delta a_\ell) = U_r U_r^\top \Delta a_\ell$, where $W = \text{col}(U_r)$.

\begin{hypothesis}[H1: Low-Dimensional Forgetting]
\label{hyp:h1}
There exists a constant $C > 0$ such that old-task performance degradation (forgetting) satisfies
\begin{equation}
    \text{forgetting} \leq C \cdot \| P_W(\Delta a) \|,
    \label{eq:h1}
\end{equation}
and the projected drift $\| P_W(\Delta a) \|$ using the top-$r$ importance-ranked dimensions is a significantly better predictor of forgetting than total drift $\| \Delta a \|$, random-dimension drift, or bottom-ranked drift.
\end{hypothesis}

This hypothesis is \emph{falsifiable}: we pre-register four predictors (top-$r$, total, random, bottom-$r$) and require the ranking top $>$ total $>$ random $>$ bottom in terms of $R^2$ for predicting forgetting across training checkpoints.

\paragraph{Ablation importance.}
Given the PCA basis $U = [u_1, \ldots, u_k]$, we define the importance of dimension $i$ as the increase in old-task loss when that dimension is zeroed out in the activation:
\begin{equation}
    \text{imp}(i) = \mathcal{L}_{\text{old}}\!\left(a_\ell - (u_i^\top a_\ell)\, u_i\right) - \mathcal{L}_{\text{old}}(a_\ell),
    \label{eq:importance}
\end{equation}
averaged over the memory set.
The top-$r$ dimensions are those with the largest importance scores.

\subsection{Sparse Feature Preservation Algorithm}
\label{sec:algorithm}

SFP operates in two phases: a \emph{basis construction} phase (before training) and a \emph{preservation} phase (during training).

\paragraph{Phase 1: Basis construction.}
Given a small memory set $\mathcal{M}$ of old-task examples (typically 128 examples), we:
\begin{enumerate}
    \item Run $\mathcal{M}$ through the pre-trained model $f_{\theta^*}$, hooking activations at layers $\ell \in \mathcal{L}$ (we use $\mathcal{L} = \{L/4, L/2, 3L/4\}$ where $L$ is the total depth).
    \item Compute the PCA basis $U_\ell \in \R^{d \times k}$ from the collected activations at each probe layer.
    \item Compute ablation importance (\cref{eq:importance}) for each PCA dimension at each layer.
    \item Select the top-$r$ dimensions per layer, yielding projection matrices $P_{W_\ell} = U_{r,\ell} U_{r,\ell}^\top$.
    \item Record reference activations $a_\ell^* = a_\ell(\theta^*)$ for the memory set.
\end{enumerate}

\paragraph{Phase 2: Preservation during training.}
The training loss combines the new-task loss with a preservation penalty:
\begin{equation}
    \mathcal{L}(\theta) = \mathcal{L}_{\text{new}}(\theta) + \lambda \sum_{\ell \in \mathcal{L}} \underbrace{\| U_{r,\ell}^\top a_\ell(\theta) - U_{r,\ell}^\top a_\ell(\theta^*) \|^2}_{\mathcal{L}_{\text{pres},\ell}(\theta)},
    \label{eq:sfp_loss}
\end{equation}
where $\lambda > 0$ controls the stability--plasticity trade-off.

The per-layer preservation loss $\mathcal{L}_{\text{pres},\ell}$ measures drift \emph{only} in the important subspace $W_\ell$.
By construction, the model is free to change activations in the complement $W_\ell^\perp$---this is where new-task learning should occur.

\paragraph{Gradient-informed variant.}
As an alternative to PCA, we also implement a gradient-informed basis: compute $G = \partial \mathcal{L}_{\text{old}} / \partial a_\ell$ on the memory set, and take the top-$r$ right singular vectors of $G$ as the preservation basis.
This directly identifies directions that the old-task loss is most sensitive to, at the cost of requiring a backward pass through the old-task loss during basis construction.

\begin{algorithm}[t]
\caption{Sparse Feature Preservation (SFP)}
\label{alg:sfp}
\begin{algorithmic}[1]
\REQUIRE Pre-trained model $f_{\theta^*}$, memory set $\mathcal{M}$, new-task data $\mathcal{D}_{\text{new}}$, probe layers $\mathcal{L}$, PCA dims $k$, preserve dims $r$, weight $\lambda$
\STATE \textbf{// Phase 1: Basis construction}
\FOR{$\ell \in \mathcal{L}$}
    \STATE Collect activations $A_\ell = \{a_\ell(\theta^*, x) : x \in \mathcal{M}\}$
    \STATE Compute PCA: $U_\ell = \text{PCA}(A_\ell, k)$
    \FOR{$i = 1, \ldots, k$}
        \STATE $\text{imp}_\ell(i) \leftarrow \E_{x \in \mathcal{M}}[\mathcal{L}_{\text{old}}(a_\ell - (u_i^\top a_\ell) u_i) - \mathcal{L}_{\text{old}}(a_\ell)]$
    \ENDFOR
    \STATE $U_{r,\ell} \leftarrow$ top-$r$ columns of $U_\ell$ by $\text{imp}_\ell$
    \STATE Store reference: $a_\ell^* \leftarrow \{a_\ell(\theta^*, x) : x \in \mathcal{M}\}$
\ENDFOR
\STATE \textbf{// Phase 2: Fine-tuning with preservation}
\FOR{each batch $B \sim \mathcal{D}_{\text{new}}$}
    \STATE $\mathcal{L} \leftarrow \mathcal{L}_{\text{new}}(B; \theta)$
    \FOR{$\ell \in \mathcal{L}$}
        \STATE $\mathcal{L} \leftarrow \mathcal{L} + \lambda \| U_{r,\ell}^\top a_\ell(\theta) - U_{r,\ell}^\top a_\ell^* \|^2$
    \ENDFOR
    \STATE Update $\theta$ via gradient descent on $\mathcal{L}$
\ENDFOR
\end{algorithmic}
\end{algorithm}

\subsection{Formally Verified Properties}
\label{sec:formal}

We formalize the mathematical properties of SFP in Lean~4 using the Mathlib library (v4.17.0).
The formalization comprises approximately 1,200 lines of Lean code with \textbf{zero} \texttt{sorry} axioms---every theorem is fully machine-checked.

\begin{definition}[Preservation Loss]
\label{def:pres_loss}
For an orthogonal projection $P_W : \R^d \to W$ onto a subspace $W$ spanned by the top-$r$ importance-ranked basis vectors, the preservation loss is
\begin{equation}
    \mathcal{L}_{\text{pres}} = \| P_W(a(\theta) - a(\theta^*)) \|^2 = \| P_W(\Delta a) \|^2.
\end{equation}
\end{definition}

The following properties are verified in Lean:

\begin{proposition}[Preservation Loss Properties, \texttt{preservationLoss\_nonneg}, \texttt{preservationLoss\_eq\_zero\_iff}]
\label{prop:loss_props}
$\mathcal{L}_{\text{pres}} \geq 0$, with equality if and only if $\Delta a \in W^\perp$ (\ie the drift is entirely in the complement of the preserved subspace).
\end{proposition}

\begin{proposition}[Pythagorean Decomposition, \texttt{drift\_pythagorean\_decomposition}]
\label{prop:pythag}
The total drift decomposes as
\begin{equation}
    \| \Delta a \|^2 = \mathcal{L}_{\text{pres}} + \| P_{W^\perp}(\Delta a) \|^2.
    \label{eq:pythagorean}
\end{equation}
Hence $\mathcal{L}_{\text{pres}} \leq \| \Delta a \|^2$ (\texttt{preservationLoss\_le\_total\_drift}): preserving the important subspace is a strictly easier constraint than preventing all drift.
\end{proposition}

\begin{proposition}[Gradient Orthogonality, \texttt{gradient\_orthogonality}, \texttt{gradient\_in\_complement}]
\label{prop:grad_orth}
At a stationary point of the combined loss $\mathcal{L}_{\text{new}} + \lambda \mathcal{L}_{\text{pres}}$ where $\mathcal{L}_{\text{pres}} = 0$, the gradient of the new-task loss satisfies $P_W(\nabla \mathcal{L}_{\text{new}}) = 0$, \ie $\nabla \mathcal{L}_{\text{new}} \in W^\perp$.
The new-task gradient is entirely in the complement of the preserved subspace.
\end{proposition}

This last property formalizes the intuition that SFP achieves a clean separation: at convergence, new-task learning occurs only in directions orthogonal to the important features.

\begin{theorem}[Main Reduction, \texttt{sfp\_reduces\_forgetting}]
\label{thm:main}
Assume H1 (\cref{hyp:h1}): $\text{forgetting} \leq C \cdot \| P_W(\Delta a) \|$ for some $C > 0$.
If training with SFP achieves $\mathcal{L}_{\text{pres}} \leq \delta$, then
\begin{equation}
    \text{forgetting} \leq C \sqrt{\delta}.
    \label{eq:reduction}
\end{equation}
\end{theorem}

\begin{proof}[Proof sketch]
By definition, $\mathcal{L}_{\text{pres}} = \| P_W(\Delta a) \|^2$.
If $\mathcal{L}_{\text{pres}} \leq \delta$, then $\| P_W(\Delta a) \| \leq \sqrt{\delta}$.
By H1, $\text{forgetting} \leq C \| P_W(\Delta a) \| \leq C\sqrt{\delta}$.
The full proof in Lean verifies that $P_W$ is an orthogonal projection (idempotent, self-adjoint) and that the norm inequality follows from the projection's contractive property (\texttt{orthogonalProjection\_norm\_le}).
\end{proof}

\begin{theorem}[Multi-Layer Extension, \texttt{sfp\_reduces\_forgetting\_multilayer}]
\label{thm:multilayer}
If forgetting decomposes across layers as $\text{forgetting} \leq \sum_\ell C_\ell \| P_{W_\ell}(\Delta a_\ell) \|$ and $\sum_\ell \mathcal{L}_{\text{pres},\ell} \leq \delta$, then $\text{forgetting} \leq (\sum_\ell C_\ell) \sqrt{\delta / |\mathcal{L}|}$ by Cauchy--Schwarz.
\end{theorem}

\begin{theorem}[Stability--Plasticity Trade-off, \texttt{stability\_plasticity\_tradeoff}]
\label{thm:tradeoff}
The preservation loss provides an upper bound on forgetting (stability) while the Pythagorean decomposition (\cref{eq:pythagorean}) guarantees that $\| P_{W^\perp}(\Delta a) \|^2 = \| \Delta a \|^2 - \mathcal{L}_{\text{pres}}$ remains available for new-task learning (plasticity).
As $\lambda \to \infty$, $\mathcal{L}_{\text{pres}} \to 0$ and forgetting is minimized, but learning is confined to $W^\perp$.
\end{theorem}

\paragraph{What is and is not verified.}
The formalization verifies the \emph{mathematical} chain from H1 to the forgetting bound.
It does \emph{not} verify H1 itself---that is an empirical claim tested in \cref{sec:experiments}.
This separation is deliberate: the formal verification makes the logical structure transparent, so the reader knows exactly which claim requires empirical evidence.
